\documentclass[11pt,a4paper]{article}
\usepackage[margin=2.5cm]{geometry}
\usepackage{amsmath}
\usepackage{booktabs}
\title{Single-cell model with transient inhibition\\\large Equations for the Brian2 tutorial}
\author{}
\date{}
\begin{document}
\maketitle

\section*{Model overview}
The model extends the spontaneously firing Hodgkin--Huxley cell with a
rectangular, conductance-based inhibitory synaptic event. The membrane equation
is
\begin{align}
C_m\frac{dV}{dt} &= I_{\mathrm{Na}}+I_{\mathrm{K}}+I_{\mathrm{L}}+I_{\mathrm{inh}},\\
I_{\mathrm{Na}} &= \bar g_{\mathrm{Na}}m^3h(E_{\mathrm{Na}}-V),\\
I_{\mathrm{K}} &= \bar g_{\mathrm{K}}n^4(E_{\mathrm{K}}-V),\\
I_{\mathrm{L}} &= g_{\mathrm{L}}(E_{\mathrm{L}}-V),\\
I_{\mathrm{inh}} &= g_{\mathrm{inh}}(t)(E_{\mathrm{inh}}-V).
\end{align}
Positive current is defined as depolarizing. Since $E_{\mathrm{inh}}=-80$ mV,
the synaptic current is normally negative and hyperpolarizing.

\section*{Synaptic event}
Let $t_{\mathrm{on}}$ be the onset and $\Delta t$ the duration. The inhibitory
conductance is a rectangular pulse,
\begin{equation}
g_{\mathrm{inh}}(t)=
\begin{cases}
\bar g_{\mathrm{inh}}, & t_{\mathrm{on}}\le t<t_{\mathrm{on}}+\Delta t,\\
0, & \text{otherwise}.
\end{cases}
\end{equation}
The tutorial uses $t_{\mathrm{on}}=0.5$ s, $\Delta t=0.05$ s, and
$\bar g_{\mathrm{inh}}=2$ mS/cm$^2$. These values are student-editable.

\section*{Gating dynamics}
For each gate $x\in\{m,h,n\}$,
\begin{equation}
\frac{dx}{dt}=\alpha_x(V)(1-x)-\beta_x(V)x.
\end{equation}
With $V$ in mV and rates in ms$^{-1}$,
\begin{align}
\alpha_m&=\frac{(V+40)/10}{1-e^{-(V+40)/10}}, &
\beta_m&=4e^{-(V+65)/18},\\
\alpha_h&=0.07e^{-(V+65)/20}, &
\beta_h&=\frac{1}{1+e^{-(V+35)/10}},\\
\alpha_n&=0.1\frac{(V+55)/10}{1-e^{-(V+55)/10}}, &
\beta_n&=0.125e^{-(V+65)/80}.
\end{align}

\section*{Parameters}
\begin{center}
\begin{tabular}{lll}
\toprule
Parameter & Value & Meaning\\
\midrule
$C_m$ & $1\,\mu$F/cm$^2$ & capacitance density\\
$\bar g_{\mathrm{Na}}$ & $120$ mS/cm$^2$ & maximal sodium conductance\\
$\bar g_{\mathrm{K}}$ & $36$ mS/cm$^2$ & maximal potassium conductance\\
$g_{\mathrm{L}}$ & $0.3$ mS/cm$^2$ & leak conductance\\
$\bar g_{\mathrm{inh}}$ & $2$ mS/cm$^2$ & inhibitory event strength\\
$E_{\mathrm{Na}}$ & $50$ mV & sodium reversal potential\\
$E_{\mathrm{K}}$ & $-77$ mV & potassium reversal potential\\
$E_{\mathrm{L}}$ & $-30$ mV & depolarizing leak reversal\\
$E_{\mathrm{inh}}$ & $-80$ mV & inhibitory reversal potential\\
\bottomrule
\end{tabular}
\end{center}

The cell starts at $V(0)=-65$ mV, with each gate initialized to
$x_\infty(V)=\alpha_x/(\alpha_x+\beta_x)$. Spikes are measured at upward
crossings of the operational sodium-activation threshold $-55$ mV. Since
Hodgkin--Huxley activation is continuous, this is a pedagogical measurement
threshold rather than a hard channel switch; voltage is not reset.
\end{document}
